Un solénoïde de \qty{\v{L}}{cm} de long est formé de \num{\v{N}} spires. Il est parcouru par un courant de \qty{\v{I}}{A}. La perméabilité du vide est $\mu_0 = \qty{1.257e-6}{T.m/A}$.

\begin{enumerate}[label=\alph*)]
    \item Déterminez l'intensité du champ magnétique au centre du solénoïde, en millitesla.
    \item On comprime le solénoïde pour que sa longueur soit réduite de moitié, sans changer le nombre de spires ni le courant. Quel est le nouveau champ magnétique ?
    \item Revenu à sa longueur initiale, quel courant faudrait-il pour obtenir un champ de \qty{\v{Bc}}{mT} ?
\end{enumerate}
